% Chapter 3 -- Scattering by a central potential; Chapter 8 -- Hamiltonian.
%
% Reference convention: Goldstein, Poole & Safko, "Classical Mechanics,"
% 3rd ed., Ch. 3 Sec. 3.10 (scattering) and Ch. 8 (Hamilton's equations),
% adopted for terminology, symbol choice and equation style. Every deviation
% from Goldstein is stated in the document itself, in the notation table of
% Sec. 3 -- a convention a reader cannot see is not a convention.

\documentclass[11pt]{article}
\usepackage{goldsteinnotes}
\usepackage{float}

\graphicspath{{figures/}}
% A table caption sits above its rules. skip=6pt was measured at ~2 pt in the
% tagged build (the setting is not honoured once tagging rewrites captions),
% so the gap is forced with \abovetopsep inside the tabular, which tagging
% does not swallow. Target: at least 6 pt of clear paper between descenders
% and the rule, verified in pixels after the build, not by the setting.
\captionsetup[table]{skip=12pt,position=top}
\captionsetup{width=\linewidth}
\setlength{\abovetopsep}{10pt}

% --- symbols fixed for THIS document -------------------------------------
% m            : mass of the incident particle
% \ell         : conserved angular momentum
% s            : impact parameter (Goldstein's own symbol; kept, not Taylor's b)
% \theta       : scattering angle -- the final, measured deflection
% \psi         : apsidal angle
% \orbang      : the running orbit angle, \orbang' its integration constant
% \varepsilon  : orbit eccentricity
% \sigma_T     : total cross section
% q_i, p_i     : generalized coordinates and conjugate momenta
\newcommand{\orbang}{\tilde\theta}
\newcommand{\Om}{\Omega}
% The orbit angle's time derivative is written as an explicit quotient rather
% than with a dot: a dot over a tilded symbol stacks two accents on one letter
% and produces a composite glyph that appears in no physics text.
\newcommand{\orbrate}{\frac{\d\orbang}{\d t}}
\newcommand{\ee}{\mathrm{e}}   % exponential, kept upright so it cannot be
                               % read as the elementary charge e

% Source errata and editorial additions are opposite claims -- one says the
% source is wrong, the other says the source left a gap -- so they get
% different devices and are never confusable.
% Star form omits the closing sentence, for the one erratum (E9) where the
% text is deliberately *not* reproduced as written.
\makeatletter
\newcommand{\erratum}{\@ifstar\erratum@nosay\erratum@say}
\newcommand{\erratum@say}[1]{\par\smallskip\noindent\emph{Erratum.} #1
  \ Reproduced above as written.\par\smallskip}
\newcommand{\erratum@nosay}[1]{\par\smallskip\noindent\emph{Erratum.} #1\par\smallskip}
\makeatother
\newcommand{\supplied}[1]{\par\smallskip\noindent\emph{Supplied.} #1\par\smallskip}

\title{Scattering by a Central Potential;\\ the Hamiltonian}
\author{HC 6001 \quad--\quad Dr.\ Castillo}
\date{22 October 2024}

\begin{document}
\maketitle

\section{Scattering by a central potential}

For a particle of mass $m$, the incident energy, the incident velocity and the
angular momentum are
%
\begin{align}
  E    &= \frac{m}{2}\,v_0^{2}, \label{eq:E}\\
  v_0  &= \sqrt{2E/m}, \label{eq:v0}\\
  \ell &= s\,m\,v_0 = s\sqrt{2mE}, \label{eq:ell}
\end{align}
%
where $s$ is the impact parameter, so that
%
\begin{equation}\label{eq:s-of-ell}
  s = \frac{\ell}{\sqrt{2mE}}.
\end{equation}
%
The geometry these quantities describe is set out in Fig.~\ref{fig:beam-ring}.

\begin{figure}
  \centering
  \includegraphics[width=0.92\textwidth]{scat_fig1_beam_ring.png}
  \caption{The scattering geometry that fixes the cross section. A particle
  approaches parallel to the axis at impact parameter $s$, turns at the
  apsidal point of closest approach, and leaves along an asymptote making the
  scattering angle $\theta$ with the forward direction; the pair of closely
  spaced straight lines from the center bracket the interval from $\theta$ to
  $\theta+\d\theta$. Particles entering through the ring of area
  $2\pi s\,\d s$ are precisely those emerging into that interval, and matching
  the two counts is what gives Eq.~\eqref{eq:sigma-general}. The particle's
  path is the curve drawn in blue and the axis and bracketing lines are black,
  a distinction of colour alone that the surrounding text states in words. The
  distance of closest approach is marked $r_<$ in the drawing and written
  $r_m$ from Eq.~\eqref{eq:psi-quadrature} onward; they are one quantity.}
  \label{fig:beam-ring}
\end{figure}

The volume element and the solid angle are
%
\begin{equation}\label{eq:solid-angle}
  \d^{3}r = r^{2}\,\d r\,
    \underbrace{\d\varphi\,\sin\theta\,\d\theta}_{\textstyle =\ \d\Om},
  \qquad
  \d\Om = 2\pi \sin\theta\,\d\theta ,
\end{equation}
%
the azimuthal angle $\varphi$ integrating to $2\pi$ because it is irrelevant by
symmetry. The quantity to be computed is $s(\theta,E)$.

\subsection{Incident intensity and differential cross section}

The incident intensity and the differential cross section are defined as
ratios of counts:
%
\begin{equation}\label{eq:intensity}
  I = \frac{\text{number of particles incident per unit time}}
           {\text{incident beam cross-sectional area}},
\end{equation}
%
\begin{equation}\label{eq:dcs-def}
  \sigma(\Om)\,\d\Om
    = \frac{\text{number of particles scattered into } \d\Om \text{ per unit time}}
           {\text{incident intensity}} .
\end{equation}
%
For a central potential $\sigma$ depends on $\Om$ only through $\theta$, so
$\sigma(\Om)$ and $\sigma(\theta)$ denote the same function throughout.

Particles crossing the incident ring between $s$ and $s+\d s$ are exactly those
emerging with scattering angle between $\theta$ and $\theta+\d\theta$
(Fig.~\ref{fig:beam-ring}), so the two counts are equal:
%
\begin{equation}\label{eq:count-match}
  I\,(2\pi s)\,\lvert \d s\rvert
    = \genfrac{}{}{0pt}{}{\text{particles through the ring}}{\text{between } s \text{ and } s+\d s}
    = \genfrac{}{}{0pt}{}{\text{particles emerging between}}{\theta \text{ and } \theta+\d\theta}
    = I\,\sigma(\theta)\,2\pi \sin\theta\,\lvert \d\theta\rvert .
\end{equation}
%
Therefore
%
\begin{equation}\label{eq:sigma-general}
  \boxeq[sigma of theta equals s over sine theta, times the absolute value of
  d s by d theta]{%
    \sigma(\theta) = \frac{s}{\sin\theta}\left\lvert \dd{s}{\theta} \right\rvert }
\end{equation}
%
Define also the total cross section
%
\begin{equation}\label{eq:sigma-T-def}
  \boxeq[sigma sub T equals the integral of d Omega times sigma of Omega]{%
    \sigma_T = \int \d\Om\,\sigma(\Om) }
\end{equation}

\subsection{The deflection angle and the apsidal angle}

The two signs of the force give the two geometries of
Figs.~\ref{fig:repulsive} and~\ref{fig:attractive}.

\begin{figure}
  \centering
  \includegraphics[width=0.80\textwidth]{scat_fig2_repulsive.png}
  \caption{Repulsive deflection. The path turns away from the center at the
  apsidal point, and the two apsidal angles $\psi$ are measured from the
  apsidal line, one back to the incoming asymptote and one forward to the
  outgoing one, which is what makes Eq.~\eqref{eq:repulsive-angles} a
  statement about a straight line.}
  \label{fig:repulsive}
\end{figure}

\begin{figure}
  \centering
  \includegraphics[width=0.72\textwidth]{scat_fig3_attractive.png}
  \caption{Attractive deflection. The path is pulled across the axis rather
  than away from it, which reverses how $\theta$ is built from the same
  apsidal angle $\psi$ and accounts for the sign difference between
  Eqs.~\eqref{eq:repulsive-angles} and~\eqref{eq:attractive-angles}.}
  \label{fig:attractive}
\end{figure}

In the repulsive case
%
\begin{equation}\label{eq:repulsive-angles}
  \pi = \theta + 2\psi,
  \qquad \theta = \pi - 2\psi,
  \qquad \psi = \frac{\pi-\theta}{2},
\end{equation}
%
and in the attractive case
%
\begin{equation}\label{eq:attractive-angles}
  \pi = 2\psi - \theta,
  \qquad \theta = 2\psi - \pi,
  \qquad \psi = \frac{\pi+\theta}{2} .
\end{equation}
%
In all cases
%
\begin{equation}\label{eq:theta-abs}
  \boxeq[theta equals the absolute value of pi minus 2 psi]{%
    \theta = \lvert \pi - 2\psi \rvert }
\end{equation}

\subsection{Computing \texorpdfstring{$s(\theta,E)$}{s(theta,E)}: the universal method}

The radial energy equation and the angular momentum together reduce the orbit
to a quadrature. Writing $\orbang$ for the running orbit angle, to keep it
distinct from the scattering angle $\theta$,
%
\begin{align}
  E &= \frac{m}{2}\dot r^{2} + V(r) + \frac{\ell^{2}}{2mr^{2}},
      \label{eq:energy-radial}\\
  \dot r &= \left[\frac{2}{m}\left(E - V(r) - \frac{\ell^{2}}{2mr^{2}}\right)\right]^{1/2},
      \qquad \dot r = \dd{r}{t} \implies \d t = \frac{\d r}{\dot r},
      \label{eq:rdot}
\end{align}
%
and, since $\ell = mr^{2}\,\orbrate$,
%
\begin{equation}\label{eq:dtheta}
  \orbrate = \frac{\ell}{mr^{2}}
  \implies
  \d\orbang = \frac{\ell}{mr^{2}}\,\d t
            = \frac{\ell}{mr^{2}}\,\frac{\d r}
              {\left[\frac{2}{m}\left(E - V(r) - \frac{\ell^{2}}{2mr^{2}}\right)\right]^{1/2}} .
\end{equation}
%
Taking $2E/m$ out of the bracket turns this into
%
\begin{equation}\label{eq:dtheta-factored}
  \d\orbang =
    \underbrace{\frac{\ell}{m}\sqrt{\frac{m}{2E}}}_{\textstyle =\ \ell/\sqrt{2mE}\ =\ s}
    \ \frac{r^{-2}\,\d r}
      {\left[1 - \dfrac{V(r)}{E}
        - \underbrace{\frac{\ell^{2}}{2mE}}_{\textstyle =\ s^{2}}\ \frac{1}{r^{2}}\right]^{1/2}} ,
\end{equation}
%
so that, integrating from the apsis out to infinity,
%
\begin{equation}\label{eq:psi-quadrature}
  \psi = \int_{0}^{\psi}\d\psi
       = \int_{\orbang_m}^{\pi}\d\orbang
       = \int_{r_m}^{\infty}
         \frac{s\,r^{-2}\,\d r}
              {\left[1 - V(r) - \dfrac{s^{2}}{r^{2}}\right]^{1/2}} ,
\end{equation}
%
where $\orbang_m$ is the orbit angle at the apsis and $r_m$ the corresponding
radius.
%
\erratum{The potential term in Eq.~\eqref{eq:psi-quadrature} is written $V(r)$
where $V(r)/E$ is meant. Eq.~\eqref{eq:dtheta-factored}, which it comes from,
carries the $/E$, and so does the result it leads to,
Eq.~\eqref{eq:theta-of-s}; without it the bracket adds an energy to two
dimensionless terms.}
%
With the substitution
%
\begin{equation}\label{eq:u-sub}
  u = \frac{1}{r}, \qquad \d u = r^{-2}\,\d r, \qquad u_m = \frac{1}{r_m},
\end{equation}
%
\erratum{$\d u = r^{-2}\,\d r$ is written without its minus sign. The sign is
absorbed exactly by the reversal of limits that turns
$\int_{r_m}^{\infty}$ into the integral from $0$ to $u_m$ below.}
%
Eq.~\eqref{eq:psi-quadrature} then becomes
%
\begin{equation}\label{eq:theta-of-s}
  \boxeq[theta of s equals the absolute value of pi minus twice the integral
  from zero to u sub m of s d u, over the square root of one minus V of u over
  E minus s squared u squared]{%
    \theta(s) = \left\lvert\, \pi - 2\int_{0}^{u_m}
      \frac{s\,\d u}{\left[1 - \frac{V(u)}{E} - s^{2}u^{2}\right]^{1/2}} \,\right\rvert }
\end{equation}
%
From here, compute
%
\begin{equation}\label{eq:dsdtheta-inv}
  \left\lvert\dd{s}{\theta}\right\rvert
    = \frac{1}{\left\lvert\dfrac{\d\theta}{\d s}\right\rvert}
\end{equation}
%
and use it to get
$\sigma(\theta) = \frac{s}{\sin\theta}\left\lvert\dd{s}{\theta}\right\rvert$,
which is Eq.~\eqref{eq:sigma-general}.

Up to here the results are valid for \emph{any} central potential $V(r)$.

\section{Coulomb scattering}

Charges are written in Gaussian units throughout, as in the reference
convention, so that the Coulomb force carries no factor of $4\pi\epsilon_0$.
The two like charges of Fig.~\ref{fig:charges} drive each other apart.

\begin{figure}
  \centering
  \includegraphics[width=0.46\textwidth]{scat_fig0_charges.png}
  \caption{The two charges $Ze$ and $Z'e$ driving each other apart along the
  line joining them, the case in which the force of
  Eq.~\eqref{eq:coulomb-force} is repulsive.}
  \label{fig:charges}
\end{figure}

The force and potential are
%
\begin{equation}\label{eq:coulomb-force}
  f = \frac{Z Z' e^{2}}{r^{2}}, \qquad
  f = -\frac{k}{r^{2}}, \qquad
  V = -\frac{k}{r}, \qquad
  k = -Z Z' e^{2} .
\end{equation}

For scattering, $E>0$, and the orbit is
%
\begin{equation}\label{eq:orbit}
  \frac{1}{r} = C\left[1 + \varepsilon\cos(\orbang - \orbang')\right],
  \qquad
  C = \frac{mk}{\ell^{2}} = -\frac{m}{\ell^{2}}\,Z Z' e^{2} ,
\end{equation}
%
with eccentricity
%
\begin{equation}\label{eq:eccentricity}
  \varepsilon = \left[1 + \frac{2E\ell^{2}}{mk^{2}}\right]^{1/2}
              = \left[1 + \frac{2E\ell^{2}}{m\,(Z Z' e^{2})^{2}}\right]^{1/2}
              = \left[1 + \left(\frac{2sE}{Z Z' e^{2}}\right)^{2}\right]^{1/2} > 1 ,
\end{equation}
%
the last step following from $\ell = s\,m\,v_0 = s\sqrt{2mE}$, which gives
%
\begin{equation}\label{eq:2l2E}
  \frac{2\ell^{2}E}{m}
    = \frac{s^{2}\,2^{2}\,\cancelled[m cancels]{m}\,E^{2}}{\cancelled[m cancels]{m}}
    = (2sE)^{2} .
\end{equation}

The orbit types this potential admits are listed in Table~\ref{tab:conics}.

\begin{table}[H]
  \centering
  \caption{Orbit type by sign of the energy. These three cases hold only for
  the potential $V=-k/r$.}
  \label{tab:conics}
  \headerrows{1}
  \begin{tabular}{ll}
    \noalign{\tablesummary{Gives the orbit type -- hyperbolic, parabolic or elliptic -- for each sign of the energy, for the potential minus k over r.}}%
    \toprule
    Energy & Orbit \\
    \midrule
    $E>0$ & hyperbolic \\
    $E=0$ & parabolic  \\
    $E<0$ & elliptic   \\
    \bottomrule
  \end{tabular}
\end{table}

There are two situations:
%
\begin{enumerate}
  \item repulsive case, $Z Z' > 0 \implies k = -Z Z' e^{2} < 0$;
  \item attractive case, $Z Z' < 0 \implies k = -Z Z' e^{2} > 0$
        (also gravity, $k = G m_1 m_2$).
\end{enumerate}

\subsection{Repulsive case}

Here $k<0$ and $C = mk/\ell^{2} < 0$, so, dividing by the negative $C$ and
reversing the inequality at each step,
%
\begin{align}
  \frac{1}{r} = C\left[1 + \varepsilon\cos(\orbang-\orbang')\right] > 0
  &\iff 1 + \varepsilon\cos(\orbang-\orbang') < 0 \notag\\
  &\iff \varepsilon\cos(\orbang-\orbang') < -1 \notag\\
  &\iff \cos(\orbang-\orbang') < -\frac{1}{\varepsilon} .
  \label{eq:repulsive-ineq}
\end{align}
%
The cosine must be negative and larger in magnitude than $1/\varepsilon$. In
this case $\orbang = \orbang'$ is not a physical point, because
$\cos(\orbang'-\orbang') = \cos 0 = 1$ would give
$\frac{1}{r} = C\left[1+\varepsilon\right] < 0$, which is unphysical. So
$\orbang = \orbang'$ lies in the middle of the inaccessible region of the
hyperbola, and $\orbang = \orbang' \pm \pi$ is the periapsis; choosing
$\orbang' = \pi$ gives
%
\begin{equation}\label{eq:cos-shift}
  \cos(\orbang-\orbang') = \cos(\orbang-\pi)
    = \cos\orbang\,\underbrace{\cos\pi}_{-1}
    + \sin\orbang\,\underbrace{\sin\pi}_{0}
    = -\cos\orbang ,
\end{equation}
%
and therefore
%
\begin{equation}\label{eq:repulsive-orbit}
  \frac{1}{r} = \left(-\frac{mk}{\ell^{2}}\right)
                \left[-1 - \varepsilon\cos(\orbang-\orbang')\right]
              = \frac{m\,Z Z' e^{2}}{\ell^{2}}\left[\varepsilon\cos\orbang - 1\right].
\end{equation}
%
This is a maximum at $\orbang = 0$, so the minimum radius is
%
\begin{equation}\label{eq:rmin-repulsive}
  \frac{1}{r_m} = \frac{m\,Z Z' e^{2}}{\ell^{2}}\left[\varepsilon - 1\right],
\end{equation}
%
and the angle at the apsidal point is $\orbang = 0$. From here on the origin of
$\orbang$ is the apsidal line, so the apsis is at $\orbang = 0$. In
Eq.~\eqref{eq:psi-quadrature} the origin was the forward beam axis, with the
incoming asymptote at $\orbang = \pi$.

What is needed now is $\psi = \orbang(r\to\infty) = \orbang(1/r \to 0)$. In the
repulsive case $\theta = \pi - 2\psi$, so $\psi = (\pi-\theta)/2$ and
%
\begin{equation}\label{eq:psi-repulsive}
  0 = \varepsilon\cos\psi - 1
  \iff
  \frac{1}{\varepsilon} = \cos\psi = \cos\!\left(\frac{\pi-\theta}{2}\right)
    = \underbrace{\cos\frac{\pi}{2}}_{0}\cos\frac{\theta}{2}
    + \underbrace{\sin\frac{\pi}{2}}_{1}\sin\frac{\theta}{2}
    = \sin\frac{\theta}{2} .
\end{equation}

\subsection{Attractive case}

Here $k>0$ and $C = km/\ell^{2} > 0$. Take $\orbang' = 0$; then
$\orbang = \orbang' = 0$ gives
%
\begin{equation}\label{eq:attractive-orbit}
  \frac{1}{r} = C\left[1 + \varepsilon\cos(0)\right] = C\left[1+\varepsilon\right],
\end{equation}
%
a maximum at $\orbang = 0$, so the minimum radius is
$1/r_m = C\left[1+\varepsilon\right]$ and the apsidal point is also at
$\orbang = 0$.

Again what is needed is $\psi = \orbang(r\to\infty) = \orbang(1/r\to 0)$. In the
attractive case $\theta = 2\psi - \pi$, so $\psi = (\pi+\theta)/2$ and
%
\begin{equation}\label{eq:psi-attractive}
  0 = 1 + \varepsilon\cos\psi
  \iff
  -\frac{1}{\varepsilon} = \cos\psi = \cos\!\left(\frac{\pi+\theta}{2}\right)
    = \underbrace{\cos\frac{\pi}{2}}_{0}\cos\frac{\theta}{2}
    - \underbrace{\sin\frac{\pi}{2}}_{1}\sin\frac{\theta}{2}
    = -\sin\frac{\theta}{2} .
\end{equation}

\section{The Rutherford cross section}

Cancelling the overall minus sign in Eq.~\eqref{eq:psi-attractive} gives
$1/\varepsilon = \sin(\theta/2)$, exactly as
Eq.~\eqref{eq:psi-repulsive} does, so both cases can be carried forward
together.

\supplied{The step joining the two cases is supplied. Each case is carried to
its own result and the two are then used together without it.}

\begin{equation}\label{eq:cot2}
  \cot^{2}\frac{\theta}{2}
    = \frac{1}{\tan^{2}\frac{\theta}{2}}
    = \frac{\cos^{2}\frac{\theta}{2}}{\sin^{2}\frac{\theta}{2}}
    = \frac{1-\sin^{2}\frac{\theta}{2}}{\sin^{2}\frac{\theta}{2}}
    = \frac{1}{\sin^{2}\frac{\theta}{2}} - 1
    = \varepsilon^{2} - 1
    = \left(\frac{2sE}{Z Z' e^{2}}\right)^{2},
\end{equation}
%
so that
%
\begin{equation}\label{eq:s-of-theta}
  \boxeq[s equals the absolute value of Z Z prime e squared over 2 E, times
  cotangent of theta over 2]{%
    s = s(\theta,E) = \left\lvert \frac{Z Z' e^{2}}{2E} \right\rvert
        \cot\frac{\theta}{2} }
\end{equation}

Differentiating,
%
\begin{align}
  \dd{s}{\theta}
    &= \left\lvert\frac{Z Z' e^{2}}{2E}\right\rvert
       \dd{}{\theta}\!\left(\frac{\cos\frac{\theta}{2}}{\sin\frac{\theta}{2}}\right)
     = \left\lvert\frac{Z Z' e^{2}}{2E}\right\rvert
       \left(\frac{1}{2}\frac{-\sin\frac{\theta}{2}}{\sin\frac{\theta}{2}}
             - \frac{1}{2}\frac{\cos^{2}\frac{\theta}{2}}{\sin^{-2}\frac{\theta}{2}}\right)
       \notag\\
    &= -\left\lvert\frac{Z Z' e^{2}}{4E}\right\rvert
        \frac{1}{\sin^{2}\frac{\theta}{2}} ,
  \label{eq:dsdtheta}
\end{align}
%
\erratum{The second term inside the bracket of Eq.~\eqref{eq:dsdtheta} is
written over $\sin^{-2}(\theta/2)$ where $\sin^{2}(\theta/2)$ is meant; the
quotient rule gives $-\tfrac{1}{2}\cot^{2}(\theta/2)$ for it, which with the
first term makes $-\tfrac{1}{2}\csc^{2}(\theta/2)$, and that is what the second
line uses.}
%
Substituting into Eq.~\eqref{eq:sigma-general} gives
%
\begin{equation}\label{eq:sigma-build}
  \sigma(\theta) = \frac{s}{\sin\theta}\left\lvert\dd{s}{\theta}\right\rvert
    = \frac{1}{2}\left(\frac{Z Z' e^{2}}{2E}\right)^{2}
      \cdot \frac{1}{2\sin\frac{\theta}{2}\,\cancelled[cosine of theta over two cancels]{\cos\frac{\theta}{2}}}
      \cdot \frac{\cancelled[cosine of theta over two cancels]{\cos\frac{\theta}{2}}}{\sin\frac{\theta}{2}}
      \cdot \frac{1}{\sin^{2}\frac{\theta}{2}} ,
\end{equation}
%
which is the Rutherford scattering cross section for Coulomb forces:
%
\begin{equation}\label{eq:rutherford}
  \boxeq[sigma of theta equals one quarter times, Z Z prime e squared over 2 E,
  all squared, times cosecant to the fourth power of theta over 2]{%
    \sigma(\theta) = \frac{1}{4}\left(\frac{Z Z' e^{2}}{2E}\right)^{2}
      \csc^{4}\!\left(\frac{\theta}{2}\right) }
\end{equation}

\subsection[Remarks on Eq. (37)]{Remarks on Eq.~\eqref{eq:rutherford}}

\begin{enumerate}
  \item The cross section is the same for the repulsive and attractive cases,
        since it is proportional to $(Z Z')^{2}$.
  \item It diverges like $E^{-2}$ at low energies.
  \item It diverges like $\theta^{-4}$ (very strongly) for small angles $\theta$.
  \item This classical result is exactly the same as the result obtained in
        quantum mechanics.
\end{enumerate}
%
\erratum*{$(zz')^{2}$ is written in lower case, $E^{-2}$ and $\theta^{-4}$ as
unrendered markup, and the third remark as ``VERY STRONG!!!''. The mathematics
is set here as mathematics and the casing brought into line with the $Z Z'$
used throughout.}

\subsection{Total scattering cross section}

The total cross section of Eq.~\eqref{eq:sigma-T-def} may now be evaluated:
%
\begin{equation}\label{eq:sigmaT}
  \sigma_T = \int \sigma(\Om)\,\d\Om .
\end{equation}
%
In this case $\sigma(\theta) = \text{const}\cdot 1/\sin^{4}(\theta/2)$, and
with
%
\begin{equation}\label{eq:dOmega-half}
  \d\Om = 2\pi\sin\theta\,\d\theta
        = 4\pi\sin\frac{\theta}{2}\cos\frac{\theta}{2}\,\d\theta ,
\end{equation}
%
\begin{equation}\label{eq:sigmaT-integral}
  \sigma_T = \frac{4\pi}{4}\left(\frac{Z Z' e^{2}}{2E}\right)^{2}
    \int_{0}^{\pi}
    \frac{\sin\frac{\theta}{2}\cos\frac{\theta}{2}}{\sin^{4}\frac{\theta}{2}}\,\d\theta ,
\end{equation}
%
the range $0<\theta<\pi$ being the full span of deflections available
(Fig.~\ref{fig:scatcenter}). The small-angle behaviour of the integrand
controls the result:
%
\begin{equation}\label{eq:divergence}
  \sigma_T \sim \int_{0}^{\delta}\frac{\d\theta}{(\theta/2)^{3}}
     = 8\int_{0}^{\delta}\frac{\d\theta}{\theta^{3}}
     = 8\left.\frac{\theta^{-2}}{(-2)}\right|_{0}^{\delta}
     = +\infty ,
\end{equation}
%
which is strongly divergent.

\begin{figure}
  \centering
  \includegraphics[width=0.60\textwidth]{scat_fig4_scattering_center.png}
  \caption{The full range of deflection directions available to a particle
  passing the scattering center, $0<\theta<\pi$, which is the range integrated
  over in Eq.~\eqref{eq:sigmaT-integral}.}
  \label{fig:scatcenter}
\end{figure}

The divergence behaves differently in the two theories, as
Table~\ref{tab:classical-quantum} records.

\begin{table}[H]
  \centering
  \caption{Behaviour of the total cross section for long-range potentials in
  the two theories.}
  \label{tab:classical-quantum}
  \headerrows{1}
  \begin{tabular}{l>{\raggedright\arraybackslash}p{0.62\linewidth}}
    \noalign{\tablesummary{Compares the classical and quantum total cross section for long-range potentials, stating for which potentials each diverges.}}%
    \toprule
    Theory & Total cross section \\
    \midrule
    Classical & $\sigma_T = +\infty$ for any long-range potential \\
    Quantum   & $\sigma_T = +\infty$ for some long-range potentials, those
                whose $V(r)$ does not decay fast enough as $r\to\infty$ \\
    \bottomrule
  \end{tabular}
\end{table}

\subsection{Notation}

Table~\ref{tab:notation} sets the symbols used here beside those of Goldstein,
Poole \& Safko, \emph{Classical Mechanics}, 3rd~ed.\ (Ch.~3.10 and Ch.~8), and
records where this document departs from that convention.

\begin{table}[H]
  \centering
  \caption{Symbols used here, their counterparts in Goldstein, Poole \& Safko,
  and the departures from that convention.}
  \label{tab:notation}
  \headerrows{1}
  \begin{tabular}{ll>{\raggedright\arraybackslash}p{0.46\linewidth}}
    \noalign{\tablesummary{Lists each symbol used here beside the symbol used in Goldstein, Poole and Safko for the same quantity, and notes where this document departs from that convention.}}%
    \toprule
    Used here & Goldstein / elsewhere & Quantity, and any departure \\
    \midrule
    $\sigma(\Om)$ & $\dd{\sigma}{\Om}$ & differential scattering cross section \\
    $\sigma_T$    & $\sigma$           & total scattering cross section \\
    $\theta$      & $\Theta$           & scattering angle; Goldstein sets it
                                         capital, and lower case is kept here \\
    $\orbang$     & $\theta$           & running orbit angle, tilded so that it
                                         cannot be read as the scattering angle \\
    $\varepsilon$ & $e$                & orbit eccentricity, kept distinct from
                                         the elementary charge $e$ \\
    $s$           & $b$                & impact parameter \\
    \bottomrule
  \end{tabular}
\end{table}

\begin{caveat}
In a dense conducting material the potential is screened,
%
\begin{equation}\label{eq:yukawa}
  V = -\frac{k}{r} \longrightarrow
  V^{\mathrm{screen}}(r) = -\frac{k}{r}\,\ee^{-r/\lambda},
\end{equation}
%
the screened Coulomb (``Yukawa'') potential. In non-conducting materials
$V^{\mathrm{screen}}(r)$ may take a different form, but will still vanish much
faster than $1/r$.
\end{caveat}

\section{Chapter 8. The Hamiltonian and Hamilton's equations of motion}

Starting from
%
\begin{equation}\label{eq:lagrangian}
  L = T - V = L(q_1,\dots,q_n,\dot q_1,\dots,\dot q_n,t),
\end{equation}
%
the Lagrange equations and the conjugate momenta are
%
\begin{equation}\label{eq:lagrange-eq}
  0 = \ddt{}\left(\evalat{\pd{L}{\dot q_i}}{q_j,\dot q_j}\right)
      - \evalat{\pd{L}{q_i}}{\dot q_j,q_j},
  \qquad
  p_i = \evalat{\pd{L}{\dot q_i}}{q_j,\dot q_j} .
\end{equation}

\begin{remark}[Midterm remark]
In polar coordinates $\partial L/\partial\dot\theta = p_\theta = \ell_z$,
where $\theta$ is the polar coordinate rather than the scattering angle of the
preceding sections. What follows from it is struck out:
$\formulaalt{retracted: L equals L of r, r-dot, p-theta}%
\xcancel{L = L(r,\dot r,p_\theta)}$,
\sout{giving the Lagrange equation}. The conserved momentum may not be
substituted back into the Lagrangian and the Lagrange equations then reused.
\end{remark}

\subsection{The Legendre transformation}

The transformation to be made is
%
\begin{equation}\label{eq:legendre-goal}
  q_1,\dots,q_n,\dot q_1,\dots,\dot q_n \longrightarrow q_1,\dots,q_n,p_1,\dots,p_n ,
\end{equation}
%
and for a function of two variables it rests on
%
\begin{equation}\label{eq:df}
  f(x,y) \implies \d f = \pd{f}{x}\,\d x + \pd{f}{y}\,\d y .
\end{equation}

\subsubsection{The thermodynamic case}

In this subsection $T$, $V$ and $p$ denote temperature, volume and pressure,
not the kinetic energy, potential and momenta of the preceding sections. The
first law, for a fluid at constant $N$, is
%
\begin{equation}\label{eq:first-law}
  \d U = \dbar Q + \dbar W = T\,\d S - p\,\d V,
  \qquad U = U(S,V) \ \text{the internal energy}.
\end{equation}
%
To have $T$ and $V$ as the variables instead, introduce the Helmholtz free
energy $F = A = U - TS$:
%
\begin{equation}\label{eq:helmholtz}
  \d A = \d U - \d(TS)
       = \cancelled[T d S cancels]{T\,\d S} - p\,\d V
         - \cancelled[T d S cancels]{T\,\d S} - S\,\d T
       = -S\,\d T - p\,\d V
  \implies A = A(T,V) .
\end{equation}

\subsubsection{The general case}

The same transformation in the abstract, with $u$ here the variable conjugate
to $x$ rather than the $1/r$ of Eq.~\eqref{eq:u-sub}, reads
%
\begin{equation}\label{eq:general-legendre}
  f(x,y) = u\,\d x + v\,\d y,
  \qquad u = \pd{f}{x}, \qquad v = \pd{f}{y},
  \qquad (x,y) \longrightarrow (u,y) .
\end{equation}
%
\erratum{The left-hand side of Eq.~\eqref{eq:general-legendre} is written
$f(x,y)$ where the differential $\d f$ is meant, as Eq.~\eqref{eq:df} has it
and as Eq.~\eqref{eq:g} then consumes it.}
%
\begin{equation}\label{eq:g}
  g = f - ux
  \implies
  \d g = \d f - \d(ux)
       = \cancelled[u d x cancels]{u\,\d x} + v\,\d y
         - \cancelled[u d x cancels]{u\,\d x} - x\,\d u
       = -x\,\d u + v\,\d y,
  \qquad g = g(u,y) .
\end{equation}

\subsection{The Hamiltonian}

Applying this to $L(q_1,\dots,q_n,\dot q_1,\dots,\dot q_n)$ with
%
\begin{equation}\label{eq:p-def}
  \ovaleq[p sub i equals the partial derivative of L with respect to q dot sub
  i, at fixed q j and q dot j]{p_i = \evalat{\pd{L}{\dot q_i}}{q_j,\dot q_j}}
\end{equation}
%
gives
%
\begin{equation}\label{eq:H-def}
  H = -\left(L - \sum_i p_i \dot q_i\right),
\end{equation}
%
written this way to mirror $g = f - ux$ of Eq.~\eqref{eq:g}, with $H = -g$; it
is the usual $\sum_i p_i\dot q_i - L$. Then
%
\begin{equation}\label{eq:dH-expand}
  \d H = -\d L + \sum_i \d(p_i \dot q_i)
       = -\pd{L}{t} -\sum_i\left(\pd{L}{q_i}\,\d q_i
         + \cancelled[the partial L by partial q dot i term cancels]
           {\pd{L}{\dot q_i}\,\d\dot q_i}\right)
         + \sum_i\left(\cancelled[the p i d q dot i term cancels]{p_i\,\d\dot q_i}
           + \dot q_i\,\d p_i\right) ,
\end{equation}
%
so that
%
\begin{equation}\label{eq:dH}
  \d H = -\pd{L}{t} + \sum_i\left(\dot q_i\,\d p_i - \pd{L}{q_i}\,\d q_i\right)
  \implies H = H(p_1,\dots,p_n,q_1,\dots,q_n,t) .
\end{equation}
%
\erratum{The first term of Eqs.~\eqref{eq:dH-expand} and~\eqref{eq:dH} carries
no differential: with every other term a differential, it is
$-\pd{L}{t}\,\d t$.}

Reading off the coefficients gives the canonical equations of motion, the
middle one using the Lagrange equation~\eqref{eq:lagrange-eq}:
%
\begin{equation}\label{eq:canonical}
  \boxeq[partial H by partial p i equals q dot i; partial H by partial q i
  equals minus partial L by partial q i equals minus p dot i; partial H by
  partial t equals minus partial L by partial t]{%
    \pd{H}{p_i} = \dot q_i,
    \qquad
    \pd{H}{q_i} = -\pd{L}{q_i} = -\dot p_i,
    \qquad
    \pd{H}{t} = -\pd{L}{t} }
\end{equation}

\end{document}
